\section{The Flow-Level Result}
\label{sec:simulation}

One instrumented vehicle cannot demonstrate throughput or delay. The flow-level
claim therefore comes from simulation. The question the simulation answers is narrow. 
It is whether restricting the
controller's observation to the fields a traffic API returns costs throughput.

\noindent\textbf{The two observations being compared.}
The prior work controller~\cite{self_regulating_cars} reads six quantities per super-segment: density, lane
count, mean speed, mean gap, entry rate and exit rate. Four of the six are
vehicle-level ground truth in the published implementation, and a vehicle cannot
obtain them. The gap is the mean of each vehicle's follow distance. The entry
and exit rates are set differences over the vehicle identities present one
minute apart. The density is counted over the vehicles on the segment. The
deployed observation of Section~\ref{sec:vision} replaces those six with the
five fields a feed and a map supply. The reward is not changed by the
substitution. Training is centralized, so it may use privileged information, and
it keeps the published form on ground-truth density. Only the state the
controller reads at decision time differs between the two arms.

On the road those five fields come from the HERE Traffic
API~\cite{here_traffic_api}, whose \emph{jam factor} is a congestion index on a
0 to 10 scale. The simulator has no such feed. Speed and free flow are read from
the road itself, and the jam factor is computed as
$\mathrm{clip}(10\,(1 - \mathrm{speed}/\mathrm{freeFlow}), 0, 10)$. HERE does
not publish the formula behind its own jam factor, and it folds in incident data
this project does not have. That expression is therefore a stated model rather than a
measurement, and it is the one place where the simulated observation and the
deployed observation are not the same quantity.

\noindent\textbf{The port.}
The published controller~\cite{self_regulating_cars} ran on the Mainz network in PTV Vissim under a
Wiedemann-99 car-following model. That network was converted to
SUMO~\cite{lopez2018sumo} directly from the Vissim XML. The conversion produces
410 edges over 51.5\,km. Three conversion decisions matter. The source network is
flat in the dynamics Vissim runs, because every one of its 410 links carries a
gradient of zero. Its geometry nevertheless carries vertical offsets of up to
six metres, which separate overpasses from the roads beneath them rather than
describing terrain. The importer reads gradients off that geometry and produces
edges far steeper than any road, up to 923\%, and SUMO applies gradient to
acceleration. Flattening the import therefore restores the source's own
dynamics. Every edge imports with a 50\,km/h limit while the action
set tops out at 60\,km/h, so a speed factor of 1.2 lifts the limit to
60\,km/h and all three actions become expressible. The exit is modelled as a three-into-two lane drop,
because the source link ends into free outflow and the network otherwise has no
bottleneck at that point. The conversion also does not import right-of-way: the
port carries 2 junctions with conflicting movements against 10 active conflict
areas in the source. All three decisions apply identically to every experiment.

\noindent\textbf{The fleet, and why it is not the published one.}
The fleet is the Extended Intelligent Driver Model with a one-second reaction
time~\cite{salles2020eidm,treiber2000idm}. It was chosen because it has a
capacity drop. A capacity drop is not a free parameter of a simulation. Real
queue discharge runs 5 to 20\% below free-flow capacity, which is a property of
traffic rather than of a model~\cite{hall1991capacity}. Measured on this port,
the Wiedemann-99 calibration discharges 4.3\% faster than free flow, which is
inconsistent with real world measurements. Its queues discharge at 1{,}980\,veh/h against 1{,}899\,veh/h
free-flowing. The Extended Intelligent Driver Model with a one-second reaction
drops 21\% at an isolated lane drop. A simulation that fails to reproduce a
measured phenomenon is a settings problem, not evidence that the simulation is
the reference. The gate applied before any experiment was run is that served flow must
fall as offered demand rises, by more than the seed spread.

\noindent\textbf{Protocol.}
Seeds 1 to 10 train, seeds 11 to 15 select the checkpoint, and seeds 16 to 30
evaluate. Four arms across fifteen seeds is 60 episodes. Each episode runs
2{,}500\,s at a 0.5\,s step, with the measurement window opening at 900\,s,
under a sustained demand of 4{,}500\,veh/h and 3{,}124 scheduled departures, at
100\% penetration.

\noindent\textbf{Why the window opens at 900\,s.}
Served flow is zero until the first trips complete, and it does not reach its
sustained level before then. Over 200\,s windows, averaged across the fifteen
evaluation seeds, the no-control arm completes 0, 43, 79 and 139 trips in the
four windows before 800\,s, and between 188 and 214 in every window after.
Measuring throughput from the start of the episode would average that fill into
the result. The network itself does not settle. Insertions exceed arrivals in
every window of the episode, so occupancy rises from 776 vehicles at the window
start to 1{,}141 at 2{,}400\,s, which is the last sample. The window makes the
throughput figure a measurement of the controller rather than of the fill, and
it does not make the episode a steady state.
Arms are paired on seed, because each arm runs the same traffic realization, and
every error bar below is two standard errors of the paired difference.

\begin{table}[t]
\centering
\caption{Held-out seeds 16 to 30 on the Mainz port, paired on seed. The
zeroed-head arm is the same network with its output layer zeroed, and it carries
the entry gate that both controlled arms carry.}
\label{tab:sim}
\small
\resizebox{\columnwidth}{!}{%
\begin{tabular}{lrrr}
\toprule
\textbf{Arm} & \textbf{Flow (veh/h)} & \textbf{Paired gain} & \textbf{Speed (km/h)} \\
\midrule
Traffic-API observation & 3{,}765 & $+230 \pm 44$ & 40.7 \\
Published six features  & 3{,}759 & $+224 \pm 61$ & 41.0 \\
No control              & 3{,}535 & reference     & 42.3 \\
Zeroed output layer     & 2{,}847 & $-688 \pm 46$ & 28.8 \\
\bottomrule
\end{tabular}
}
\end{table}

\noindent\textbf{The result.}
Table~\ref{tab:sim} gives it. Restricting the observation to what a traffic API
returns costs nothing that this experiment can resolve. The two controlled arms
differ by 6\,veh/h against error bars of 44 and 61\,veh/h, which means that
performance of the deployed version is inseparable from prior work. This shows that
our model of decentralized execution does not lose against the centralized controller.

Both controlled arms run slower than no control while serving more vehicles.
Space-mean speed falls from 42.3 to 40.7\,km/h and arrivals rise from 1{,}940 to
1{,}995 per episode. That is the intended mechanism. Holding density below
critical upstream of the merge lets the merge discharge near capacity instead of
breaking down.

\noindent\textbf{Testing the boundary.}
Both controlled arms hold a vehicle at an entry link while that link is at or
above critical density. The no-control arm does not. This gate is a boundary condition on the
controlled system rather than a piece of road. Applying it to an uncontrolled
run would credit that run with a control action it is not taking. That asymmetry is a candidate explanation for the gain, so it
was tested. The zeroed-head arm is the same trained network with its output
layer zeroed after loading. It carries the entry gate, and it loses
$688 \pm 46$\,veh/h against no control, which is 15.7 times the no-control arm's
own bar below it. The gate alone does not produce the gain.

\noindent\textbf{What this result is not comparable to.}
No number here is comparable to the published table in \cite{self_regulating_cars}. Every claim is against this
port's own no-control baseline. Table~\ref{tab:pair} states what the two papers
cover between them. The controller and the shape of the topology are what the
two runs have in common. Almost everything that decides how the controller
behaves differs, which is why the pair is stronger evidence than a matched
number in one simulator would be. A topology or demand sweep here would
re-establish what the earlier paper already did.

\begin{table}[t]
\centering
\caption{What the two papers cover between them. The road network is nominally
the same, and the port differs from the source in right-of-way, fleet and exit
geometry.}
\label{tab:pair}
\small
\begin{tabular}{lll}
\toprule
 & \textbf{Earlier paper} & \textbf{This paper} \\
\midrule
simulator & PTV Vissim & SUMO \\
car following & Wiedemann-99 & EIDM, 1\,s reaction \\
demand & 18{,}000\,veh/h surge & 4{,}500\,veh/h sustained \\
network & Mainz & Mainz \\
\bottomrule
\end{tabular}
\end{table}

\noindent\textbf{The simulated gate and the deployed gate answer different
questions.}
In simulation the advisory is written as a speed command, and SUMO bounds it by
the car-following model's own safe speed. A vehicle whose leader is slower
follows its leader. That is the same semantics the published implementation
has, and it is the documented way to command a speed in SUMO. We left it alone.
Replacing the simulator's car following with our own layer would make this arm
incomparable to the earlier paper, whose own car-following model is the row
that differs in Table~\ref{tab:pair}. It would also make the run reproducible
only by someone holding our code, rather than by someone holding SUMO and the
network file.

On the road that bound would be the wrong instrument.
SUMO's safe speed is an actuator command, and it works because the simulator
sets the vehicle's speed. The prototype sets nothing. The driver does the car
following, so the vehicle already carries the model the simulator supplies, and
a second copy of it on the phone would recompute what the driver is already
doing. The question the road poses is whether a recommendation is fit to show,
and that is what the deployed filter decides. One rule raises a recommendation
that would have the vehicle travel slowly on an uncongested road. The other
withholds the number when the gap to the vehicle ahead, or the time left to
close it, falls below the configured minimum. Neither has an analogue in the
simulator, because a simulated vehicle has no display and no driver reading it.

